% AUTO-GENERATED by the update_record tool. DO NOT EDIT.
% verify_paper regenerates this file from record.json and the run outputs
% and fails on any difference, so manual edits always surface.
\noindent\textbf{Hypothesis 1.} NVIDIA の open-weight モデル nemotron-3.5-lightning（30B、活性 3B の MoE）と nemotron-3-ultra（550B、活性 55B の MoE）を RIKYU の GPU ノード上の NIM サーバー（OpenAI 互換）で自前 serving し、SciGym 公開コードの Controller をそのまま用いて論文と同じ条件（SciGym-small 137 件、最大 20 反復、temperature 1.0）で解かせたとき、いずれも平均 Simulation Trajectory Error（STE）が論文 Table 1 の最良行（Gemini-2.5-Pro の 0.3212）以下である。~\cite[s1.p1, s1.p5, s1.p6]{duan-2025-measuring}, \cite[s3.p2]{scigymrikyuopenmodels-2026-scigym}
\begin{enumerate}
\item[\textbf{Claim 1}] nemotron-3.5-lightning（nvidia/nemotron-3.5-lightning）で SciGym-small 137 件を最大 20 反復で解かせたときの平均 STE は、論文 Table 1 の最良行（Gemini-2.5-Pro）の 0.3212 以下である。~\cite[s1.p5]{duan-2025-measuring}
  \emph{Rationale:} 活性 3B の軽量 MoE が論文の最良行に届けば、h1 の lightning の部分を担う。
  \emph{Criterion:} \texttt{\detokenize{nemotron-3.5-lightning-small}}.\texttt{\detokenize{trajectory_smape}} $-$ 0.3212 $\leq 0$.
  \emph{Prediction:} $[-0.07, 0.18]$ (同条件の 27B〜35B の open-weight モデルは 0.248〜0.457（s3.p5、s3.p6）。活性 3B の軽量モデルはこの幅に収まると見る).
  \emph{Observed:} 0.257174.
  \emph{Verdict:} refuted.
\item[\textbf{Claim 2}] nemotron-3-ultra（nvidia/nemotron-3-ultra-550b-a55b）で SciGym-small 137 件を最大 20 反復で解かせたときの平均 STE は、論文 Table 1 の最良行（Gemini-2.5-Pro）の 0.3212 以下である。~\cite[s1.p5]{duan-2025-measuring}
  \emph{Rationale:} 活性 55B の大型 MoE が論文の最良行に届けば、h1 の ultra の部分を担う。
  \emph{Criterion:} \texttt{\detokenize{nemotron-3-ultra-small}}.\texttt{\detokenize{trajectory_smape}} $-$ 0.3212 $\leq 0$.
  \emph{Prediction:} $[-0.22, -0.02]$ (同条件の大型 open-weight MoE（kimi-k2.6、kimi-k3、deepseek-v4.1-flash）は 0.21〜0.24（s3.p2〜s3.p4）。ultra は同等かやや良いと見る).
  \emph{Observed:} -0.0376069.
  \emph{Verdict:} supported.
\end{enumerate}
\noindent\emph{Assumed, so that the claims together imply Hypothesis 1:}
\begin{itemize}
\item temperature 1.0 での 1 回実行の平均 STE（137 件平均）の実行間変動が判定に対して十分小さいこと。先行する Table 1 再現では件ごとの STE の標準偏差はおよそ 0.27、137 件平均の標準誤差はおよそ 0.02 であった（c1、c2）
\item NIM コンテナ（vLLM）で serving したモデルが、提供元の推奨設定で呼んだ場合と同じ能力で応答すること。思考（reasoning）はモデルの既定のままとし、思考の文が本文に含まれて返る場合もそのまま Controller に渡す（c1、c2）
\item aarch64 向けの libroadrunner 2.7.0 と antimony 2.14.0 が、論文環境の libroadrunner 2.8.0 と同じ軌道と評価値を与えること（c1、c2）
\item 平均 STE が「論文の最良行に届くか」を代表すること。RMS（modifier あり / なし）と NTS は同じ実行から表として報告するが、判定基準には用いない（c1、c2）
\item 文脈長（262144 トークン）や時間切れで提出に至らなかった件は不完全モデルの採点値で平均に含める（c1、c2）
\item NIM サーバーの job が計算機の 4 日上限で打ち切られても、別ノードで立ち直したサーバーに実験側が接続し直し、反復の途中で条件が変わらないこと（c1、c2）
\end{itemize}
\noindent\textbf{Hypothesis 2.} 同じ 2 モデルを、論文が評価していない SciGym-large 213 件（反応数の多い系）で同じ条件（最大 20 反復、temperature 1.0）で解かせたとき、いずれも平均 STE が論文 Table 1 の最良行（small での Gemini-2.5-Pro の 0.3212）以下である。~\cite[s1.p2, s1.p5]{duan-2025-measuring}, \cite[s3.p2]{scigymrikyuopenmodels-2026-scigym}
\begin{enumerate}
\item[\textbf{Claim 3}] nemotron-3.5-lightning で SciGym-large 213 件を最大 20 反復で解かせたときの平均 STE は、論文 Table 1 の最良行（Gemini-2.5-Pro）の 0.3212 以下である。~\cite[s1.p5]{duan-2025-measuring}
  \emph{Rationale:} 軽量モデルが大規模な系でも論文の最良行に届けば、h2 の lightning の部分を担う。
  \emph{Criterion:} \texttt{\detokenize{nemotron-3.5-lightning-large}}.\texttt{\detokenize{trajectory_smape}} $-$ 0.3212 $\leq 0$.
  \emph{Prediction:} $[0.03, 0.28]$ (large には先行値が無い。small での予測（0.25〜0.50）に、反応数の多い系での悪化を 0.1 程度見込む).
  \emph{Observed:} 0.275656.
  \emph{Verdict:} refuted.
\item[\textbf{Claim 4}] nemotron-3-ultra で SciGym-large 213 件を最大 20 反復で解かせたときの平均 STE は、論文 Table 1 の最良行（Gemini-2.5-Pro）の 0.3212 以下である。~\cite[s1.p5]{duan-2025-measuring}
  \emph{Rationale:} 大型モデルが大規模な系でも論文の最良行に届けば、h2 の ultra の部分を担う。
  \emph{Criterion:} \texttt{\detokenize{nemotron-3-ultra-large}}.\texttt{\detokenize{trajectory_smape}} $-$ 0.3212 $\leq 0$.
  \emph{Prediction:} $[-0.12, 0.13]$ (large には先行値が無い。small での予測（0.10〜0.30）に、反応数の多い系での悪化を 0.1 程度見込む).
  \emph{Observed:} 0.111574.
  \emph{Verdict:} refuted.
\end{enumerate}
\noindent\emph{Assumed, so that the claims together imply Hypothesis 2:}
\begin{itemize}
\item large split には論文の公表値が無いので、論文の small での最良値を絶対水準の基準とする。small との難易度の差は表で並べて示すが、判定には用いない（c3、c4）
\item 213 件平均の実行間変動が判定に対して十分小さいこと。件ごとの標準偏差を 0.27 とすれば平均の標準誤差はおよそ 0.02（c3、c4）
\item 評価層 scigym\_large（airas-eval v0.13.1）が公式リリースの large split 213 件を sha256 で固定し、small と同じ公式 Evaluator で採点すること（c3、c4）
\item 文脈長（262144 トークン）や時間切れで提出に至らなかった件は不完全モデルの採点値で平均に含める（c3、c4）
\end{itemize}
